% 30 Mar: clean up

\subsection{Installment 2: “Folkebladet,” March 20, 1918}

When I took up Sverdrup’s Collected Writings, the principles, and so forth, I studied them without any preconceived opinion. I studied them in the light of the experiences I had had, and the knowledge I had gathered.

When I read the second volume of the Collected Writings, the volume that bears the title The Congregation, I had to say: this volume has for the most part to do with church polity.

What a spiritually dead conclusion! I hear someone say. Can Eyen not see that the volume The Congregation deals directly with awakening and life?

But no, I classified the book under ecclesiastical literature on church polity. Nor was I the only one who did so. Pastor S. E. Sørgensen,\footnote{Pastor S. E. Sørgensen, born in 1846, cand. theol. 1873, missionary pastor in the service of the Norwegian Missionary Society in Madagascar, where he directed the Norwegian Mission’s seminary, 1875–1886. Director of the Mission School in Stavanger, 1892–97. Later pastor in Trondheim. Illustreret Norsk Konversations Leksikon says: “S. has been one of the most prolific authors in the field of missionary literature. We mention here Missionslære [Theory of Missions], Trold og Kirke paa Madagaskar [Sorcery and the Church in Madagascar], Vidnesbyrd fra Missionen [Testimonies from the Mission], Missionsprækener [Mission Sermons], and many articles concerning church questions and missions in church periodicals and the daily press. In the Malagasy language he has written a larger church history, pastoral theology, liturgics, and several other commentaries.} of Trondhjem, Norway, did the same, when in 1910 he wrote:

\begin{quote}
“The second volume consists almost exclusively of articles from the long debate over questions of church polity that has been carried on among our countrymen in America.”
\end{quote}

And he added: “The rich and mighty Reformed church life pulsating all around him (Sverdrup) would seem to have been able to furnish him with more than enough analogies, at any rate. But it is remarkable how small a role all this plays, if indeed it can be said to be touched upon at all.”

Thus, Sørgensen held that the second volume of Sverdrup’s Collected Writings was directed toward church polity, and that the Reformed church life could have furnished Sverdrup with more than enough analogies. Unfortunately, however, we were told by our pastors that there had never been any discussion of church polity among the Norwegian Lutherans in America! Reflections force themselves upon one.

Was Sørgensen mistaken? Was it lack of spirit, of experience, of theological formation, that caused him to write thus?

I do not think so. I believe he could hold his own with any of us in these matters. As a competent theologian he gave the matter its right name.

Sørgensen spoke, as we see, of analogies. I have more than once called attention to these analogies in my lectures on church polity. And I hope this will not be regarded as some unpardonable offense.

What Sørgensen says of Volume II may, with greater or lesser reservation, also be said of Volumes III and IV.

By comparing Sverdrup’s articles on the congregation with the polity literature that stands at my disposal, I find exceedingly little in Sverdrup for which I do not find analogies in my works on church polity, which fill a shelf six feet long. I speak perhaps speaking as a fool before the wise, but others do the same also, for example Harvard’s five-foot shelf of classics. I do not at all mean to say that Sverdrup drew from this literature. But no one can deny the kinship: it is intimate.

Now there are doubtless those who would dismiss this six-foot shelf as a mere little list, and its contents as trivial stuff. But there is something strange about this content: it is counted trivial on that shelf, but becomes at once spirit and life when it comes between the covers of the Collected Writings. One wonders whether opponents are justified in drawing the distinction they have habitually drawn: polity versus spirit and life.

What is true of most of Sverdrup’s and Oftedal’s contributions to church debate holds especially for “Guiding Principles.” All these principles — they are eight, in twelve paragraphs — are polity principles, neither more nor less.

Any recognized Protestant theological faculty in Scandinavia or Germany that is allowed to see the principles will be able to tell us that. That they are of significance for spirit and life, no one will deny. But one may say that about all work for the kingdom of God, yes, also about collections of money for seminary, mission, and so forth. Week after week the receipts lay claim to many columns in Folkebladet. It never comes to an end, yet no one complains that they are not “spiritual.” At times they represent no little spirit.

Guiding Principles are principles of church polity. Pastor Caspersen has come into a tremendously mystical vein when in certain articles in Folkebladet he maintains that the question ‘Do you want a congregation?’ and ‘Guiding Principles’ concern the person as such, or that sets forth a “new type of Christian.” It is almost incomprehensible that a man with theological formation can write such a thing.

Quite otherwise thought, wrote, and spoke Pastor E. Lou, who knew his Bible as few among us did, and who kept up with theology and labored zealously as perhaps none of us did. Lou wrote an immense amount in former days in Luthersk Kirkeblad and in Folkebladet. He both thought and theologized, but in an altogether different, sober, unpretentious manner. In June 1914, when he reviewed Guide to the Free Church’s Principles, he wrote in Folkebladet:


\begin{quote}
If the people of the Free Church do not agree on this, that the principles are matters of church polity, then the Free Church’s days are numbered.
\end{quote}

Words almost prophetic, whose fulfillment perhaps lies nearer than many think.

And this was not Lou’s view only in 1914. In the last letter he wrote to me, shortly before his death — and I have permission to quote him — he said: “We agree that the principles are matters of church polity and not an exhortation to ‘awakening’ in this word’s meaning, as the ‘awakened’ understand it.”

Now it may well be rather indifferent what we call the principles. If we had called them practical theology, that name would doubtless have given less offense than the more precise term principles of church polity, for church polity is, after all, a part of practical theology.

But since exceedingly different understandings of the principles prevail with respect to what they mean — and as to their value there will always be difference of opinion — it is surely right that they be made the object of discussion until agreement is reached on what they actually entail. For the Free Church gains little by having principles that mean one thing to some, another thing to others. It is meaningless to connect them with a “new type of Christian,” and it is likewise meaningless to confuse the Free Church’s task or purpose with its principles. The key to their understanding lies in the literature of church polity.

Pastor Holck in Norway, one of the friends of reform, who in his day wrote much about questions of church polity, wrote in the 1870s: “Indirectly it has been stated that the supreme principle for the church must be to look to its own welfare, to its purpose, to bring souls to salvation, and that to this concern every other consideration must give way. In spite of its appearance of truth, this rests upon a manifest misunderstanding. To advance the church’s purpose, the salvation of souls, can never become a principle in the same sense in which religious liberty is one; it is a task that is to be solved.”

When I hear that the principles are time and again classified under “spirit and life” and similar slogans, I see a picture: a schoolman lays four open books on a table, all in the English language; the one book treats astronomy, the second physiology, the third physics, and the fourth philosophy. “This is astronomy,” says the schoolman, “this is physiology, this is physics, this is philosophy.” No and again no, comes the answer from a stubborn onlooker: “It is neither the one nor the other: it is all English!” There is no making anything precise for such a fellow.\footnote{Historically offensive content removed.}

But if we could agree that the principles concern church polity, would that mean that the Free Church’s chief task is to labor for church polity?

No, every church’s foremost task must be the salvation of souls. But if that were all that mattered, then one could simply refrain from organizing congregations, let alone associations of congregations. There must be something more.

In the Free Church’s Rules — rules are not the same as principles — we read that the purpose of the association called the Lutheran Free Church is to labor “for the quickening and liberation of all Norwegian Lutheran congregations.” Is it only the awakening and salvation of souls that are here being aimed at?

Those who authored the principles should here be allowed to speak. Sverdrup wrote: “It is not merely a question of awakening the individual conscience. It is not solely the task to bring the individual to the Savior. The task is to awaken consciousness of the congregation’s essence and calling, by the power of the Word and the Spirit, to build up and people the body of Christ.” (Collected Writings III.)

And in other places he writes of the manifestation of the body of Christ, that the congregation shall win its proper position among us, and so forth.

It is the duty of all Christians to labor for the quickening and liberation of individuals. But it is not thereby said that this, without more ado, means “the quickening and liberation of the congregations.” In Norway much Christian work is done — I am inclined to think that there is even more fear of God among the church people there than among us. But in spite of this Sverdrup would not admit that in the Norwegian state church there are congregations in the proper sense of the word. I do not state my own opinion here, but Sverdrup’s — and precisely in order to shed light on the principles. After a visit to Norway in 1886 Sverdrup wrote this:

\begin{quote}
The state church has too little room for Christian and ecclesiastical work.... There is no congregation in the proper sense of the word in which there is room for all gifts and use for all hands. If some Christian man or woman feels a call and inclination to some work, then they must break their own way on their own. There is no one who marks out the field of labor for them....

Naturally there is room for many works of love and much self-sacrifice in family life, in the school, in private intercourse with people. But the actual specifically Christian field of labor is lacking, the congregation, the body of Jesus Christ. In all outward respects it is so bound by law, and in a certain sense so fatherly provided for by the state, that there is no room for the Christians’ self-sacrificing labor for it.

Christian people who come from Norway to America often find that the brethren in the Free Church are so dry and cold, so filled with busyness, that it is almost doubtful whether they can really have any true heart Christianity, any inward life in God. They seem far too outwardly oriented and absorbed in work.

Conversely, he who comes from America to Norway gets the impression that Christianity is chiefly feeling, inwardness, contemplation, thought, but little or no work....

But it would be altogether unjust to conclude from this that there is therefore less true and honest fear of God in the one place than in the other.”
\end{quote}

Thus: the fear of God is to be equally true and honest in both lands. And yet: no congregation in the proper sense of the word in Norway.

How then can the sort of language we hear so much of be grounded in Sverdrup, language which says that if only souls are converted and brought to God, then one is thereby, simply, laboring for the quickening and liberation of the congregation?

Certainly, one must begin with individuals, but one cannot remain standing with individuals. If there should be found a little congregation in which all were believers, could one then say that it did not need to be counted among those congregations to which the Free Church refers when in its rules it speaks of “the vivification and liberation of all Norwegian Lutheran congregations”?

No, according to the Augsburg Confession’s understanding of the congregation, something more is required: there must be labor so that the congregation becomes a visible congregation, which, because it bears the name of the invisible congregation, must correspond to that name.

Therefore it is not consistent to say that the organizations are indifferent. But it is tempting to speak thus, which Sverdrup also emphasizes when he says:

\begin{quote}
In the New Testament such high and holy words are spoken about the congregation that many have given it up in despair and have sought to escape from the holy demand by teaching that all the great and glorious things are said about an invisible congregation, and that is — another matter.”

The body of Christ, that whereby Christ is seen, that through which Christ works, is visible, because it is believing men and women who labor with the Word and the Sacraments, and it is precisely these outwardly organized fellowships that are called congregations.” (1897.)
\end{quote}

Here the Norwegian Lutheran Free Church distinguishes itself from the “societies” and the “people’s churches” and the “state churches,” which make the congregations into wholly subordinate members of the larger organizations or at any rate do not set them higher than these organizations. The Free Church regards the local congregation as a divine institution; it regards the other organizations as human arrangements.

Pastor P. L. Tangjerd has in Four Hundred Years of Lutheranism criticized the Free Church for this. He says that sentences such as “the congregation is of divine origin, but the association is of human origin,” “the local congregation is the proper form of the kingdom of God on earth” — are “suited to confuse the concepts.”

He continues:

\begin{quote}
“What is a congregation? Disciples gathered around the means of grace (‘two or three in Jesus’ name’). Whether they have an outward organization, how long the gathering lasts, how near one another the disciples live, and so forth, makes no difference. An annual meeting of the church body is, in the biblical sense, just as much a congregation as any local congregation within it. Both are equally much and equally little ‘instituted by God’ and ‘the proper form of the kingdom of God on earth’.... ‘Congregation’ and ‘association’ stand side by side with respect to divine institution....
\end{quote}
